\section{Methodology}

\subsection{Overview}

We compare OLMo-7B~\cite{Groeneveld2024OLMo} and Pythia-6.9B~\cite{Biderman2023Pythia} at matched training scale using publicly available checkpoints. Both models are evaluated on MMLU~\cite{Hendrycks2021MMLU}, HellaSwag~\cite{Zellers2019HellaSwag}, and ARC~\cite{Clark2018ARC} using a standardized evaluation harness. The primary metric is the MMLU/HellaSwag generalization balance ratio; the secondary metric is the ARC-Challenge minus ARC-Easy delta. Statistical robustness is assessed via bootstrap confidence intervals and effect size estimation.

\subsection{Model Selection and Training Scale Matching}

We select checkpoints that minimize deviation from 300B training tokens using each model's publicly documented checkpoint schedules.

\begin{table}[h]
\centering
\caption{Checkpoint selection for matched-scale comparison.}
\begin{tabular}{lllll}
\toprule
Model & HuggingFace ID & Revision & Approx.\ Training Tokens & Deviation \\
\midrule
Pythia-6.9B & \texttt{EleutherAI/pythia-6.9b} & \texttt{step143000} & $\approx$299.9B & $-$0.03\% \\
OLMo-7B & \texttt{allenai/OLMo-7B} & \texttt{step68000-tokens301B} & $\approx$301B & $+$0.33\% \\
\bottomrule
\end{tabular}
\end{table}

Both checkpoints are within 0.5\% of the 300B token target, providing a well-controlled matched-scale comparison. Checkpoint provenance is fully documented via HuggingFace revision hashes.

\subsection{Evaluation Protocol}

All benchmarks are evaluated using \texttt{lm-evaluation-harness} v0.4.12~\cite{gao2024harness} with default task configurations. MMLU is evaluated in 5-shot mode across all 57 subjects. HellaSwag uses 10-shot evaluation. ARC-Easy and ARC-Challenge use 25-shot evaluation. A 500-sample fast evaluation mode is used for primary results to enable rapid iteration; full benchmark evaluation is recommended to confirm HellaSwag convergence findings.

\begin{table}[h]
\centering
\caption{Evaluation benchmark configurations.}
\begin{tabular}{llll}
\toprule
Benchmark & Task Name & Shots & Metric \\
\midrule
MMLU & \texttt{mmlu} & 5 & Accuracy (macro avg.) \\
HellaSwag & \texttt{hellaswag} & 10 & Normalized accuracy \\
ARC-Easy & \texttt{arc\_easy} & 25 & Accuracy \\
ARC-Challenge & \texttt{arc\_challenge} & 25 & Accuracy \\
\bottomrule
\end{tabular}
\end{table}

\subsection{Primary and Secondary Metrics}

The \textbf{primary metric} is the MMLU/HellaSwag generalization balance ratio:

\begin{equation}
    R = \frac{\text{MMLU}_{\text{acc}}}{\text{HellaSwag}_{\text{acc}}}
\end{equation}

This ratio captures the balance between knowledge acquisition (MMLU numerator) and commonsense generalization (HellaSwag denominator). A higher ratio indicates stronger relative knowledge acquisition given commonsense baseline performance.

The \textbf{secondary metric} is the ARC difficulty delta:

\begin{equation}
    \Delta_{\text{ARC}} = \text{ARC-Challenge}_{\text{acc}} - \text{ARC-Easy}_{\text{acc}}
\end{equation}

A more negative delta indicates stronger relative degradation on harder reasoning questions, which may reflect differences in reasoning robustness between corpora.

\subsection{Statistical Analysis}

We use non-parametric bootstrap resampling ($B = 1000$ iterations) to construct 95\% confidence intervals for the ratio difference (OLMo $-$ Pythia). At each bootstrap iteration, we resample subjects (for MMLU) or examples (for HellaSwag) with replacement, compute per-model ratios, and record the difference. Cohen's $d$ is computed using the pooled standard deviation of per-subject MMLU accuracy as the effect size measure. The null hypothesis is $H_0: R_{\text{OLMo}} \geq R_{\text{Pythia}}$; rejection requires the 95\% CI to lie entirely below zero.

\subsection{Limitations of the Comparison Design}

This comparison is not architecture-controlled. OLMo-7B and Pythia-6.9B differ in attention mechanism, tokenizer vocabulary, learning rate schedule, and optimizer configuration. Performance differences are therefore a joint function of corpus quality and architecture, and causal attribution to corpus quality alone is not warranted. We discuss this limitation in detail in Section~6 and characterize the conditions under which the architecture confound would need to be ruled out to support corpus quality conclusions.
